\documentclass[12pt,a4paper]{article}
\usepackage[utf8]{inputenc}
\usepackage[T1]{fontenc}
\usepackage[french]{babel}
\usepackage{geometry}
\usepackage{fancyhdr}
\usepackage{amsmath}
\usepackage{listings}
\usepackage{xcolor}
\usepackage{hyperref}
\usepackage{graphicx}
\usepackage{enumitem}
\usepackage{tcolorbox}

\geometry{margin=2.5cm}
\setlength{\headheight}{14.5pt}

\pagestyle{fancy}
\fancyhf{}
\fancyhead[L]{Examen Blanc 2 - Administration Unix}
\fancyhead[R]{\thepage}
\fancyfoot[C]{}

\lstset{
    language=bash,
    basicstyle=\ttfamily\small,
    keywordstyle=\color{blue},
    commentstyle=\color{gray},
    stringstyle=\color{red},
    numbers=left,
    numberstyle=\tiny\color{gray},
    breaklines=true,
    frame=single,
    showstringspaces=false
}

\definecolor{questionbox}{RGB}{255,250,240}
\newtcolorbox{questionbox}{
    colback=questionbox,
    colframe=orange!50!black,
    title=Question,
    fonttitle=\bfseries
}

\title{\textbf{Examen Blanc 2 - Administration Unix/Linux}\\
\large Durée : 2 heures - Documents autorisés}
\author{}
\date{}

\begin{document}

\maketitle
\thispagestyle{fancy}

\section*{Instructions}
\begin{itemize}
    \item Durée : 2 heures
    \item Documents autorisés : Résumés de cours et TP
    \item Toutes les questions sont obligatoires
    \item Soignez la présentation et l'orthographe
    \item Pour les commandes, précisez si elles nécessitent des privilèges root
\end{itemize}

\newpage

% ============================================================
% QUESTION 1 : PROCÉDURE D'ADMINISTRATION
% ============================================================

\section{Question 1 : Procédure d'Administration (15 points)}

\begin{questionbox}
Vous devez configurer un nouveau serveur Linux. Décrivez étape par étape la procédure complète pour :

\begin{enumerate}
    \item Créer un utilisateur "admin" avec les privilèges sudo complets
    \item Configurer un nouveau disque de 500 Go : partitionner, formater en ext4, et monter sur /mnt/storage de manière permanente
    \item Installer et configurer le service Apache (httpd) pour qu'il démarre automatiquement au boot
    \item Créer un fichier swap de 4 Go et l'activer de manière permanente
\end{enumerate}

Pour chaque étape, donnez les commandes exactes et expliquez ce que vous faites.
\end{questionbox}

\newpage

% ============================================================
% QUESTION 2 : HIÉRARCHIE ET COMMANDES
% ============================================================

\section{Question 2 : Hiérarchie des Fichiers et Commandes (12 points)}

\begin{questionbox}
\subsection*{Partie A : Répertoires système (4 points)}
Expliquez le rôle et donnez des exemples de contenu pour :
\begin{itemize}
    \item /etc
    \item /var
    \item /usr
    \item /sbin
\end{itemize}

\subsection*{Partie B : Commandes de recherche (4 points)}
\begin{enumerate}[label=\alph*)]
    \item Comment trouver tous les fichiers .conf modifiés dans les 7 derniers jours dans /etc ?
    \item Comment rechercher le mot "error" dans tous les fichiers .log de /var/log ?
    \item Quelle commande permet de connaître le chemin complet de la commande "python3" ?
    \item Comment lister tous les fichiers de plus de 50 Mo dans /home, triés par taille décroissante ?
\end{enumerate}

\subsection*{Partie C : Liens et inodes (4 points)}
\begin{enumerate}[label=\alph*)]
    \item Créez un lien symbolique /usr/local/bin/python pointant vers /usr/bin/python3. Quelle commande utilisez-vous ?
    \item Comment vérifier qu'un fichier est un lien symbolique ?
    \item Expliquez pourquoi on ne peut pas créer de lien physique entre deux systèmes de fichiers différents.
    \item Comment connaître le nombre d'inodes utilisés sur le système de fichiers racine ?
\end{enumerate}
\end{questionbox}

\newpage

% ============================================================
% QUESTION 3 : DÉMARRAGE ET GRUB
% ============================================================

\section{Question 3 : Démarrage et Configuration GRUB (10 points)}

\begin{questionbox}
\subsection*{Partie A : Théorie (5 points)}
\begin{enumerate}[label=\alph*)]
    \item Expliquez le rôle de chaque composant dans la séquence de démarrage : BIOS/UEFI, GRUB, noyau Linux, initramfs, systemd.
    \item Pourquoi l'initramfs est-il nécessaire dans certains cas ? Donnez des exemples.
    \item Quelle est la différence entre un runlevel (SysVinit) et un target (systemd) ?
\end{enumerate}

\subsection*{Partie B : Configuration GRUB (3 points)}
Vous voulez modifier le délai d'attente du menu GRUB à 10 secondes. Décrivez la procédure complète avec les commandes.

\subsection*{Partie C : Arrêt système (2 points)}
\begin{enumerate}[label=\alph*)]
    \item Quelle commande permet d'arrêter le système de manière gracieuse ?
    \item Comment programmer un arrêt dans 1 heure avec un message pour les utilisateurs ?
\end{enumerate}
\end{questionbox}

\newpage

% ============================================================
% QUESTION 4 : GESTION DES SERVICES
% ============================================================

\section{Question 4 : Gestion des Services systemd (10 points)}

\begin{questionbox}
\subsection*{Partie A : Commandes systemd (5 points)}
Pour le service "sshd" (SSH), donnez les commandes pour :
\begin{enumerate}[label=\alph*)]
    \item Vérifier son état actuel
    \item Le démarrer s'il est arrêté
    \item L'activer pour qu'il démarre automatiquement au boot
    \item Consulter ses logs depuis le dernier démarrage
    \item Le redémarrer
\end{enumerate}

\subsection*{Partie B : Journalisation (3 points)}
\begin{enumerate}[label=\alph*)]
    \item Comment afficher les logs système depuis 1 heure ?
    \item Comment suivre les nouveaux logs en temps réel ?
    \item Comment afficher uniquement les erreurs du système ?
\end{enumerate}

\subsection*{Partie C : Targets (2 points)}
\begin{enumerate}[label=\alph*)]
    \item Comment passer du mode graphique au mode multi-utilisateurs texte ?
    \item Quelle commande permet de définir le mode multi-utilisateurs texte comme défaut au démarrage ?
\end{enumerate}
\end{questionbox}

\newpage

% ============================================================
% QUESTION 5 : GESTION DES PAQUETAGES
% ============================================================

\section{Question 5 : Gestion des Paquetages (10 points)}

\begin{questionbox}
\subsection*{Partie A : Concepts (4 points)}
\begin{enumerate}[label=\alph*)]
    \item Expliquez ce qu'est un dépôt (repository) dans le contexte de la gestion des paquets.
    \item Quelle est la différence entre "yum install" et "rpm -i" ? Quand utiliseriez-vous chacun ?
    \item Qu'est-ce qu'une dépendance transitive ? Donnez un exemple.
    \item Où sont stockées les métadonnées des paquets RPM installés ?
\end{enumerate}

\subsection*{Partie B : Commandes YUM (4 points)}
\begin{enumerate}[label=\alph*)]
    \item Comment rechercher un paquet contenant la commande "gcc" ?
    \item Comment mettre à jour tous les paquets du système ?
    \item Comment désinstaller un paquet et toutes ses dépendances non utilisées ?
    \item Comment vider le cache YUM pour libérer de l'espace disque ?
\end{enumerate}

\subsection*{Partie C : Dépôts (2 points)}
\begin{enumerate}[label=\alph*)]
    \item Où se trouvent les fichiers de configuration des dépôts YUM ?
    \item Comment activer un dépôt désactivé nommé "epel" ?
\end{enumerate}
\end{questionbox}

\newpage

% ============================================================
% QUESTION 6 : UTILISATEURS ET PERMISSIONS
% ============================================================

\section{Question 6 : Utilisateurs, Groupes et Permissions (13 points)}

\begin{questionbox}
\subsection*{Partie A : Fichiers de configuration (4 points)}
\begin{enumerate}[label=\alph*)]
    \item Décrivez la structure du fichier /etc/group. Que signifie chaque champ ?
    \item Pourquoi le fichier /etc/shadow a-t-il des permissions restrictives (400 ou 000) ?
    \item Que contient le champ "passwd" dans /etc/passwd et que signifie-t-il ?
    \item Comment modifier la date d'expiration d'un compte utilisateur ?
\end{enumerate}

\subsection*{Partie B : Gestion utilisateurs (5 points)}
\begin{enumerate}[label=\alph*)]
    \item Créez un utilisateur "developpeur" avec :
    \begin{itemize}
        \item Groupe primaire "developers"
        \item Groupes secondaires "wheel" et "docker"
        \item Shell /bin/bash
        \item Répertoire personnel /home/developpeur
    \end{itemize}
    \item Comment supprimer complètement un utilisateur et son répertoire personnel ?
    \item Comment forcer un utilisateur à changer son mot de passe à la prochaine connexion ?
\end{enumerate}

\subsection*{Partie C : Permissions et sudo (4 points)}
\begin{enumerate}[label=\alph*)]
    \item Expliquez la permission 640 en mode octal. À quoi correspond-elle en mode symbolique ?
    \item Donnez la commande pour donner les permissions rwx------ à un fichier.
    \item Comment configurer sudo pour qu'un utilisateur "backup" puisse exécuter uniquement la commande /usr/bin/tar sans mot de passe ?
    \item Quelle est la différence entre setuid et setgid sur un répertoire ?
\end{enumerate}
\end{questionbox}

\newpage

% ============================================================
% QUESTION 7 : GESTION DES DISQUES
% ============================================================

\section{Question 7 : Gestion des Disques et Systèmes de Fichiers (10 points)}

\begin{questionbox}
\subsection*{Partie A : Partitionnement (4 points)}
\begin{enumerate}[label=\alph*)]
    \item Quelles sont les principales différences entre MBR et GPT ? Dans quels cas utiliseriez-vous GPT ?
    \item Comment lister tous les disques et leurs partitions avec leurs tailles et types de systèmes de fichiers ?
    \item Quelle commande permet d'obtenir l'UUID d'une partition ? Pourquoi préférer UUID à /dev/sdX dans /etc/fstab ?
\end{enumerate}

\subsection*{Partie B : Systèmes de fichiers (3 points)}
\begin{enumerate}[label=\alph*)]
    \item Quels sont les avantages d'ext4 par rapport à ext2 ?
    \item Qu'est-ce que la journalisation dans un système de fichiers ? Quels sont ses avantages ?
    \item Donnez la commande pour formater une partition /dev/sdc1 en XFS.
\end{enumerate}

\subsection*{Partie C : Montage (3 points)}
\begin{enumerate}[label=\alph*)]
    \item Expliquez le rôle et la structure du fichier /etc/fstab.
    \item Écrivez une ligne /etc/fstab pour monter /dev/sdb1 sur /mnt/backup en ext4, en lecture seule, sans vérification fsck.
    \item Comment forcer le remontage d'un système de fichiers en lecture seule sans le démonter ?
\end{enumerate}
\end{questionbox}

\newpage

% ============================================================
% QUESTION 8 : SWAP
% ============================================================

\section{Question 8 : Gestion de l'Espace Swap (10 points)}

\begin{questionbox}
\subsection*{Partie A : Concepts (5 points)}
\begin{enumerate}[label=\alph*)]
    \item Expliquez le mécanisme de la mémoire virtuelle et le rôle du swap.
    \item Qu'est-ce que le "thrashing" ? Comment l'éviter ?
    \item Quelle est la différence entre une partition swap et un fichier swap ? Quels sont les avantages/inconvénients de chacun ?
    \item Qu'est-ce que le paramètre "swappiness" ? Comment le modifier ?
\end{enumerate}

\subsection*{Partie B : Configuration pratique (5 points)}
Décrivez la procédure complète pour créer une partition swap de 8 Go sur /dev/sdb2, l'activer, et la rendre permanente. Incluez toutes les commandes nécessaires.
\end{questionbox}

\newpage

\section*{Fin de l'examen}

\textbf{Bon courage !}

\end{document}

